% vim: ai sw=2 spell spelllang=cs
\shorthandoff{-}

\begin{frame}
  \frametitle{C a \file{libc}}

  \begin{itemize}
    \item \file{libc} se stará o komunikaci s OS (jádrem)
    \begin{itemize}
      \item Jádro poskytuje pouze ,,základní'' funkcionalitu
    \end{itemize}

    \pause

    \item \file{libc} poskytuje spoustu nadstaveb
    \begin{itemize}
      \item Soubory: \code{fopen}, \code{fread}, \code{fwrite}, \code{fclose},
	\dots

      \pause

      \item Síť: \code{getaddrinfo}, \code{ntoh*}, \code{hton*}

      \pause

      \item Vlákna (\file{libpthread}): \code{pthread_*} (včetně zamykání)

      \pause

      \item A spoustu dalšího (celý POSIX), \dots
    \end{itemize}
  \end{itemize}

  \medskip
  \pause

  \begin{center}
    \heading{Bez \file{libc} nic z toho nemáme}

    \smallskip

    Jak teď ale např.\ otevřu soubor?
  \end{center}

\end{frame}

\begin{frame}
  \frametitle{Komunikace s jádrem}

  \begin{itemize}
    \item Sdílená paměť
    \begin{itemize}
      \item Jádro zapisuje, uživatel čte a naopak
%      \item V jednom z dalších cvičení
    \end{itemize}

    \pause

    \item Systémová volání
    \begin{itemize}
      \item Podobné jako zavolání funkce
      \item Ale program a jádro běží v jiných kontextech

      \pause

      \item Přepnutí kontextu
      \begin{itemize}
	\item Uložení a změna stavu procesoru
	\item Relativně drahá operace
      \end{itemize}

      \pause

      \item Uvidíme dnes dále
    \end{itemize}

  \end{itemize}

\end{frame}

\begin{frame}
  \frametitle{Systémová volání}

  \begin{itemize}
    \item Volání funkce

    \pause

    \item Instrukcí procesoru
    \begin{itemize}
      \item Závislé na architektuře
      \item Softwarové přerušení (x86\_32: číslo \code{0x80})
      \item Speciální instrukce (x86\_32: \code{sysenter},
        x86\_64: \code{syscall})

      \pause

      \item \heading{Demo:} implementace \code{syscall} z \file{glibc}
    \end{itemize}

    \pause

    \item Systémová volání jsou očíslovaná
    \begin{itemize}
      \item Do jednoho registru číslo

      \pause

      \item Čísla \code{__NR_*} definovaná jádrem (pro každou architekturu)

      \pause

      \item \header{sys/syscall.h}
      \item Např. (x86\_64): \code{__NR_write} (1), \code{__NR_exit} (60)
    \end{itemize}
  \end{itemize}

\end{frame}

\begin{frame}
  \frametitle{Úkol}

  \heading{\file{libc} o úroveň níže}
  \begin{enumerate}
    \item Zavolejte systémová volání
    \begin{itemize}
      \item \code{write} s nějakým textem
      \item \code{exit}
    \end{itemize}

    \item Proveďte ale pomocí funkce \code{syscall} z \file{libc}
    \begin{itemize}
      \item Tedy nepište přímo \code{write(...)}, \code{exit(...)}
      \item Viz \cmd{man 2 syscall}
      \item První argument: číslo systémového volání
      \item Další argumenty: odpovídají už danému volání
      \item Např.\ \code{syscall(__NR_kill, -1, SIGKILL)}
    \end{itemize}

    \item Přeložte a vyzkoušejte
  \end{enumerate}

\end{frame}

\begin{frame}
  \frametitle{C bez \file{libc}}

  \begin{itemize}
    \item Žádné hlavičkové soubory z \file{libc}, ani jeho funkce
    \item Jen to, co poskytuje jádro (\file{/usr/include/linux/})
    \item Při překladu: \cmd{gcc -nostdinc ...}
    \item Při linkování: \cmd{gcc -nostdlib ...}
    \begin{itemize}
      \item \cmd{gcc} stále může vkládat volání \code{memset} apod.
    \end{itemize}
  \end{itemize}

\end{frame}

\begin{frame}{Volací konvence}

  \begin{center}
    \begin{tabular}{|l||c|c||c|c|} \hline
      \multicolumn{1}{|c||}{\multirow{3}{*}{Parametr}} &
        \multicolumn{4}{c|}{Registry} \\ \cline{2-5}
	 & \multicolumn{2}{c||}{Uživatelský prostor} &
	   \multicolumn{2}{c|}{\onslide<2->{Systémová volání}} \\ \cline{2-5}
	 & \xesl    & \xesh & \onslide<2->{\xesl} & \onslide<2->{\xesh} \\ \hline\hline
      1. & EAX      & RDI   & \onslide<2->{EBX}   & \onslide<2->{RDI}  \\ \hline
      2. & EDX      & RSI   & \onslide<2->{ECX}   & \onslide<2->{RSI}  \\ \hline
      3. & Zásobník & RDX   & \onslide<2->{EDX}   & \onslide<2->{RDX}  \\ \hline
      4. & Zásobník & RCX   & \onslide<2->{ESI}   & \onslide<2->{R10}  \\ \hline
      5. & Zásobník & R8    & \onslide<2->{EDI}   & \onslide<2->{R8}   \\ \hline
      6. & Zásobník & R9    & \onslide<2->{EBP}   & \onslide<2->{R9}   \\ \hline
      7. & Zásobník & Zásobník &
	\multicolumn{2}{c|}{\multirow{2}{*}{\onslide<2->{\emph{Jen 6 parametrů}}}} \\ \cline{1-3}
      8. & Zásobník & Zásobník & \multicolumn{2}{c|}{} \\ \hline\hline
      Návratová & EAX  & RAX & \onslide<2->{EAX}  & \onslide<2->{RAX}  \\ \hline
    \end{tabular}
  \end{center}

  \onslide<3->

  Pozn.: \xesl a \texttt{-mregparm}

\end{frame}

\begin{frame}
  \frametitle{Úkol}

  \heading{Systémová volání bez \file{libc}}
  \begin{enumerate}
    \item Otevřete a projděte si \file{pb173-bin/08/}
    \item Zavolejte
    \begin{itemize}
      \item \code{fork}
      \item Z potomka: \code{write} na standardní výstup
      \item Z rodiče: \code{read} 16 bytů a jejich \code{write} na standardní
	výstup
      \item Z obou potom: \code{exit}
    \end{itemize}

    \item Přeložte a spusťte
    \begin{itemize}
      \item \cmd{gcc -nostdlib ...}
    \end{itemize}
  \end{enumerate}

\end{frame}

\begin{frame}
  \frametitle{Data z jádra}

  \begin{itemize}
    \item Jádro předává
    \begin{itemize}
      \item Parametry programu
      \item Proměnné prostředí
      \item Rozšiřující vektor
    \end{itemize}

    \pause

    \item Vše na zásobníku
    \begin{itemize}
      \item Formát zásobníku je pevně daný
    \end{itemize}
  \end{itemize}

\end{frame}

\begin{frame}
  \frametitle{Formát zásobníku}

  \begin{center}
    \begin{tabular}{|l|l|} \hline
      Počet parametrů & \code{long}   \\ \hline\hline
      1.\ parametr    & \code{char *} \\ \hline
      \dots           & \code{char *} \\ \hline
      m-tý parametr   & \code{char *} \\ \hline
      Konec parametrů & \code{0UL}    \\ \hline\hline

      1.\ proměnná prostředí     & \code{char *} \\ \hline
      \dots                      & \code{char *} \\ \hline
      n-tá proměnná prostředí    & \code{char *} \\ \hline
      Konec proměnných prostředí & \code{0UL} \\ \hline\hline

      Rozšiřující vektor 1    & \code{struct \{ long type, val; \}} \\ \hline
      \dots                   & \code{struct \{ long type, val; \}} \\ \hline
      Rozšiřující vektor o    & \code{struct \{ long type, val; \}} \\ \hline
      Konec rozšiřujících vektorů & \code{\{AT_NULL, ? \}} \\ \hline

    \end{tabular}
  \end{center}

\end{frame}

\begin{frame}[fragile]
  \frametitle{Rozšiřující vektor}

  \begin{block}{Typ záznamů}
  \begin{lstlisting}
struct {
  long type;
  long value;
};
  \end{lstlisting}
  \end{block}

  \smallskip

  \begin{itemize}
    \item Typy: \code{AT_NULL}, \code{AT_ENTRY}, \code{AT_RANDOM}, \dots
    \item Pole struktur zakončuje typ \code{AT_NULL}
  \end{itemize}

\end{frame}

\begin{frame}
  \frametitle{Úkol}

  \heading{Zjištění názvu programu a platformy}
  \begin{enumerate}
    \item Pište do programu \file{pb173-bin/08/nostd.c}

    \item Vypište první parametr programu
    \begin{itemize}
      \item Tj.\ název samotného programu
    \end{itemize}

    \item Dále iterujte přes zásobník až k rozšiřujícímu vektoru
    \item Tam najděte typ \code{AT_PLATFORM}
    \item Vypište hodnotu jako řetězec
    \begin{itemize}
      \item Použijte \code{write}
    \end{itemize}

    \item Přeložte (bez \file{libc}) a spusťte
  \end{enumerate}

\end{frame}


\begin{frame}
  \frametitle{Virtuální systémová volání}

  \heading{\code{vsyscall}}
  \begin{itemize}
    \item Jen na některých architekturách
    \begin{itemize}
      \item A konfiguracích jádra
    \end{itemize}

    \item Neprobíhá přepnutí kontextu
    \item Speciální stránka(y) s kódem namapovaným jádrem
    \item Podpora jen 3 funkcí
    \begin{itemize}
      \item \code{gettimeofday}: aktuální čas (\code{0xffffffffff600000})
      \item \code{time}: čas v sekundách od Epochy (\code{0xffffffffff600400})
      \item \code{getcpu}: číslo CPU a NUMA uzlu (\code{0xffffffffff600800})
    \end{itemize}

    \item \heading{Demo:} \file{/proc/self/maps}

%    \pause
%
%    \item Adresa funkce se získá pomocí makra \code{VSYSCALL_ADDR()}
%    \begin{itemize}
%      \item \header{asm/vsyscall.h}
%      \item Parametr je název shora s předponou \code{__NR_v}
%    \end{itemize}
  \end{itemize}

\end{frame}

\begin{frame}
  \frametitle{Úkol}

  \heading{Volání \code{time} z \code{vsyscall}}
  \begin{enumerate}
    \item Zavolejte adresu s \code{time}
    \item Vypište hodnotu jako řetězec
    \item Přeložte a spusťte několikrát
  \end{enumerate}

\end{frame}

\begin{frame}
  \frametitle{Virtuální dynamická knihovna}

  \heading{\code{vdso}}
  \begin{itemize}
    \item Speciální knihovna přilinkovaná jádrem
    \item Podpora také závislá na architektuře

    \item Podobná \code{vsyscall}, ale více flexibilní
    \begin{itemize}
      \item Navíc obsahuje ale jen \code{clock_gettime}
    \end{itemize}

    \pause

    \item \heading{Demo:} \cmd{ldd} a \cmd{readelf}

    \pause

    \item Knihovna v ELF formátu
    \begin{itemize}
      \item Adresa začátku v auxv: \code{getauxval(AT_SYSINFO_EHDR)}
      \item Dále se přečte ELF pomocí \file{libelf}
      \item Ukázkový kód: \href{https://git.kernel.org/cgit/linux/kernel/git/torvalds/linux.git/tree/tools/testing/selftests/vDSO/parse_vdso.c?id=f9b6b0ef60349cf1747d8f366f23900671f888c5}{\nolinkurl{tools/testing/selftests/vDSO/parse_vdso.c}}
    \end{itemize}
  \end{itemize}

  \iffimuni
  \pause
  \bigskip

  \heading{Úkol:} domácí úkol
  \fi

\end{frame}
